\section{Los CPLD}
    \subsection{Definición}
        \sangria{} Es un array de bloques lógicos basados en suma de productos (AND-OR), interconectados por una matriz central y cuya configuración se almacena en una momemoria no volátil.
        \sangria{} Sus siglas hacen referencia a: \textit{Complex Programmable Logic Device}. Las características que lo destacan son que retienen su configuración sin alimentación (por la memoria no volátil) y el instant-on (al encender ya está listo instantáneo).

    \subsection{Fabricantes destacados y software}
        \sangria{} Es el fabricante Altera la encargada de perfeccionar los CPLD, destacando su serie EP300 que incluye una ventana de cuarzo para exponer la parte de la memoria, permitiendo su borrado usando luz ultravioleta. Junto a la serie MAX7000, contribuyeron al éxito de Altera en la época.
        \sangria{} Su programación era en HDL, concretamente Altera usaba Quartus para sus CPLD. El ''compilador'' se encargaba de mapear a la estructura AND-OR y escribirla en memoria no volátil.
